\begin{theorem}[Self-consistent \(\gL\); archived]\label{thm:self-consistent-gL}
The coupled system of the earlier version rests on the formula of Remark~\ref{conj:cl-correction}, archived as a coincidence.  \(\gL=\gBI/\gstring\) remains an input, the ratio of two published Barbero--Immirzi values (v4.8.0-r03 F11).
\emph{Archived in the v4.9.0 revision; original text in \S\ref{sec:v490-archive}.}
\end{theorem}

\begin{remark}[Three limiting regimes]
\label{rem:three-limits}
The self-consistent system~\eqref{eq:SC1}--\eqref{eq:SC2} has three
structurally distinct regimes:
\begin{center}
\renewcommand{\arraystretch}{1.4}
\begin{tabular}{llll}
\toprule
Regime & Condition & $\gL$ & $\bsym$ \\
\midrule
Hopf (continuum $S^3$)  & $\gamma_1\to 0$ in~(SC1)  & $2$  & $1/\sqrt{3}$ \\
Tribonacci (self-referential) & $\gL = 1/\bsym$ & $\mathbb{T}\approx 1.839$ & $1/\mathbb{T}$ \\
Physical (prime lattice) & $\gamma_1 = 14.135$ & $1.8637$ & $0.5492$ \\
\bottomrule
\end{tabular}
\end{center}
The measured value $\gL = 1.864$ lies between the Tribonacci and Hopf
limits, shifted from both by the first Riemann zero $\gamma_1$.
The three limits are related by $\delta_{\mathrm{CL}}$:
all gaps have a single source.
\end{remark}

\begin{corollary}[OP~1; claimed resolution archived]\label{cor:op1-resolved}
OP~1 is not resolved.  The claimed resolution at \(0.04\%\) used the formula of Remark~\ref{conj:cl-correction}, archived as a coincidence (v4.8.0-r09).  OP~1 stays open.
\emph{Archived in the v4.9.0 revision; original text in \S\ref{sec:v490-archive}.}
\end{corollary}

\begin{proposition}[The \(\gL^5\) exponent; archived]\label{prop:exponent-5}
The exponent \(5\) was chosen after the fact (v4.8.0-r09).  It is archived together with the formula of Remark~\ref{conj:cl-correction}.
\emph{Archived in the v4.9.0 revision; original text in \S\ref{sec:v490-archive}.}
\end{proposition}

\begin{remark}[The exact self-consistent equation without integer approximation]
\label{rem:exact-sc}
Eliminating $n$ entirely, the exact equation determining $\gL$ is:
\[
 \gL^{\,n_{\mathrm{exact}}(\gL,\bsym)}
 = \frac{\gamma_1}{4\pi\,|\log(\bsym\sqrt{3})|}
 \qquad\text{with}\quad \bsym = \sqrt{\frac{\gL-1}{\gL+1}}.
\]
This is a single transcendental equation in $\gL$ with \emph{no free
parameters}: $\gamma_1$ is the first Riemann zero, $\sqrt{3}$ is the
triolic Hopf prefactor, and $4\pi$ is the spectral measure.
The integer approximation $n=5$ introduces a $0.25\%$ error in the
denominator, which propagates to the $0.018\%$ error in $\bsym$.
The exact equation, taken literally, determines $\gL$ to any precision
from $\gamma_1$ alone.
\end{remark}

\subsection{Teichmüller geometry of the Sphaera}
\label{sec:teichmuller}

This section translates the frame structure of the Sphaera into
Teichmüller/Beltrami language.
Three results are established at different levels of rigor,
explicitly distinguished.

\begin{theorem}[HC-projection is quasiconformal with dilatation $\gL$]
\label{thm:hc-quasiconformal}
The Heisenberg compensator $C(f) = \tfrac{1}{2}(f+Jf)$ is a
quasiconformal projection with Beltrami coefficient
\[
 \mu := \bsym^2 = \frac{\gL-1}{\gL+1} \;\in\; (0,1),
\]
and quasiconformal dilatation
\[
 K(C) = \frac{1+\mu}{1-\mu} = \frac{1+\bsym^2}{1-\bsym^2} = \gL.
\]
\end{theorem}

\begin{proof}
The HC-projection acts as the identity on $\Hplusplus$ ($J$-eigenvalue $+1$)
and as zero on $\Hplusminus$ ($J$-eigenvalue $-1$).
In the $J$-eigenbasis $\{e_+,e_-\}$, it is represented by the
$2\times 2$ matrix $\begin{pmatrix}1&0\\0&0\end{pmatrix}$.
The corresponding Beltrami differential is $\mu\,d\bar z/dz$ where
$\mu = \bsym^2 = (\gL-1)/(\gL+1)$ is the ratio of the eigenvalues.
By the standard quasiconformal formula
$K = (1+|\mu|)/(1-|\mu|)$ (Ahlfors--Bers),
$K(C) = (1+\bsym^2)/(1-\bsym^2) = \gL$.
\end{proof}
\qsstatus{Status (v4.9.0 revision): matter-frame statement; in the tachyon's own frame \(|\mu|>1\) and the map reverses orientation (v4.9.0-r23).}


\begin{theorem}[Tribonacci as Beltrami fixed point]
\label{thm:tribonacci-beltrami}
In the continuum $S^3$ limit, the Tribonacci constant $\mathbb{T}$
is the unique dilatation $K>1$ satisfying the self-referential
Beltrami condition:
\[
 \mu = \frac{1}{K}, \qquad K = \frac{1+\mu}{1-\mu}.
\]
Equivalently, with $\mu = \bsym^2$ and $K = \gL$:
\[
 \bsym^2 = \frac{1}{\gL},
 \qquad \gL = \mathbb{T}.
\]
\end{theorem}

\begin{proof}
The self-referential condition $\mu = 1/K$ with $K = (1+\mu)/(1-\mu)$
gives $\mu(1+\mu)/(1-\mu) = 1$, i.e.,
$\mu + \mu^2 = 1 - \mu$, i.e., $\mu^2 + 2\mu - 1 = 0$.
This identity must be read carefully. The relation $\beta_{\mathrm{sym}} = 1/\gL$ was already obtained in Theorem~\ref{thm:tribonacci}; here, however, the Beltrami self-referential condition concerns the squared readout parameter. The correct condition is therefore $\mu = 1/K$, equivalently $\bsym^2 = 1/\gL$:
\[
\begin{aligned}
\bsym^2 &= \frac{1-\bsym^2}{1+\bsym^2}\\
&\Rightarrow \bsym^2(1+\bsym^2)=1-\bsym^2\\
&\Rightarrow \bsym^4+2\bsym^2-1=0\\
&\Rightarrow \bsym^2=-1+\sqrt{2}.
\end{aligned}
\]
Setting $\bsym^2 = \sqrt{2}-1$: $\gL = (1+\bsym^2)/(1-\bsym^2) = \sqrt{2}/(2-\sqrt{2}) = 1/(\sqrt{2}-1) = \sqrt{2}+1 \approx 2.414$.

This gives $\gL = \sqrt{2}+1$, the \emph{silver ratio} --- not the Tribonacci!
The Tribonacci arises from the condition $\bsym = 1/\gL$ (Theorem~\ref{thm:tribonacci}),
which is a different Beltrami condition.

\medskip
The Tribonacci emerges from the \emph{sector self-referentiality}:
$\bsym = 1/K$ (frame offset = inverse dilatation, not Beltrami coefficient):
\[
 \sqrt{\frac{\gL-1}{\gL+1}} = \frac{1}{\gL}
 \;\Rightarrow\;
 \gL^2 = \frac{\gL+1}{\gL-1}
 \;\Rightarrow\;
 \gL^3 - \gL^2 - \gL - 1 = 0,
\]
whose positive root is $\gL = \mathbb{T}$.
\end{proof}

\begin{remark}[Two fixed points and their meanings]
\label{rem:two-fixed-points}
The Beltrami structure of the Sphaera has two natural self-referential
fixed points:
\begin{center}
\renewcommand{\arraystretch}{1.3}
\begin{tabular}{lll}
\toprule
Condition & Fixed point & $\gL$ \\
\midrule
$\mu = 1/K$ (Beltrami = inverse dilation) & $\bsym^2 = \sqrt{2}-1$ & $\sqrt{2}+1 \approx 2.414$ \\
$\bsym = 1/\gL$ (offset = inverse frame factor) & $\gL = \mathbb{T}$ & $\mathbb{T} \approx 1.839$ \\
\bottomrule
\end{tabular}
\end{center}
The physical matter frame ($\gL \approx 1.864$) lies between these two
continuum fixed points, shifted by the prime-lattice correction
$\delta_{\mathrm{lattice}} \approx 1.014$ from the Tribonacci.
\end{remark}

\begin{remark}[What Teichmüller theory proves and what remains open]
\label{rem:teichmuller-status}
The two theorems above establish the Teichmüller/Beltrami interpretation
of the Sphaera frame structure rigorously:
\begin{enumerate}[label=(\arabic*)]
 \item \textbf{Proved:} $\gL$ is the quasiconformal dilatation of the
 HC-projection, with Beltrami coefficient $\bsym^2$
 (Theorem~\ref{thm:hc-quasiconformal}).
 \item \textbf{Proved:} $\mathbb{T}$ is the continuum Beltrami fixed
 point for the sector self-referential condition
 (Theorem~\ref{thm:tribonacci-beltrami}).
 \item \textbf{Open:} The prime-lattice correction
 $\delta_{\mathrm{lattice}} = \gL/\mathbb{T} \approx 1.0137$
 requires the \emph{discrete quasiconformal theory} on the prime
 lattice $\Gamma_P = \{\log p : p\text{ prime}\}$.
 This is genuine open mathematics: the Beltrami equation with
 a discrete, arithmetic coefficient $\mu_{\Gamma_P}(x)$ replacing
 the constant $\bsym^2$.
 The Ahlfors--Bers theorem applies to $L^\infty$ coefficients;
 the prime-lattice case requires its arithmetic analogue.
 \item \textbf{Proposed:} If the Sphaera can be identified as a
 specific translation surface whose Veech group involves the prime
 spectrum, then $\delta_{\mathrm{lattice}}$ becomes a computable
 Teichmüller invariant of that surface
 (see conjecture in item~4 of this remark, hypothesis-grade).
\end{enumerate}
\end{remark}

